\chapter{Capital-Structure Arbitrage}\label{ch:s2:capital-structure-arbitrage}

A company's shares and its credit default swaps both price its default risk. When they disagree, a trader buys one, sells the other, and waits for the
disagreement to close, or for the model that measured it to prove wrong. Yu tested the trade on 135,759 daily CDS spreads of 261 North American
companies, using the CreditGrades model. Single trades could lose heavily, because spreads and share prices were only loosely correlated. An equally
weighted portfolio of all trades earned Sharpe ratios like those of other fixed-income arbitrage strategies. On this chapter's sixty synthetic firms
the trade earns a Sharpe ratio of 1.40, and 48\% of its trades lose. When leveraged buyouts, which double a firm's debt while its shares jump, come
four times as often, the Sharpe ratio falls to 0.39. The build is \texttt{firm.capstruct}.

\section{From equity to credit spread}

A structural model (Book~6, chapter~14) treats equity as a claim on a firm's assets and default as the assets falling below what is owed. The version
used here is in the spirit of CreditGrades. Assets per share are the share price plus the part of the debt that would be recovered, $V = E + LD$ for
debt per share $D$ and recovered share $L$. Asset vol is the equity vol scaled down, $\sigma_E E/V$, and default is the first time $V$ touches $LD$.
The five-year survival probability $q$ from the first-passage formula gives a spread, $(1 - R)(-\ln q)/5$ for CDS recovery $R$.

\begin{definition}[Equity-implied spread]\label{def:s2:capital-structure-arbitrage:implied}
The \emph{equity-implied spread}\index{equity-implied spread} of a company is the credit default swap spread that a structural model gives from its
share price, equity volatility and debt per share; comparing it with the market spread measures whether the equity and credit markets disagree about
the company's default risk.
\end{definition}

With a share price of 50, debt of 40 a share, 35\% equity vol, $L = 0.5$ and $R = 0.4$, the implied five-year spread is 55.8 basis points
(\cref{fig:s2:capital-structure-arbitrage:curve}). At debt of 80 a share it is 115.4. Spreads rise as the share falls, but only to 189.8 basis
points at a share price of 10 with debt of 40. The model's asset vol falls with the equity, a feature of the approach, which keeps distressed spreads
lower than a constant-vol model would.

\begin{omfigure}
\begin{tikzpicture}
\begin{axis}[width=10cm,height=5.8cm,xlabel={\small share price},ylabel={\small five-year spread (bp)},tick label style={font=\footnotesize},
  xmin=10,xmax=100,ymin=0,ymax=240,grid=major,grid style={gray!25},table/col sep=comma,
  legend style={font=\scriptsize,at={(0.97,0.97)},anchor=north east}]
\addplot[omThm,thick] table[x=E,y=d40]{figdata/strategies-2/09-capital-structure-arbitrage/curve.csv};
\addplot[omDef,thick,dashed] table[x=E,y=d80]{figdata/strategies-2/09-capital-structure-arbitrage/curve.csv};
\addplot[black,thick,->,>=stealth] coordinates {(50,55.8) (60,98.8)};
\legend{debt 40 a share,debt 80 a share}
\end{axis}
\end{tikzpicture}

\omcaption{The chapter's equity-implied five-year spread by share price, at 35\% equity vol, for debt of 40 and 80 a share. The arrow is a leveraged
buyout: the share rises from 50 to 60 on the premium while the debt doubles, and the spread goes from 55.8 to 98.8 basis points. Data:
\texttt{s2\_capstruct.curve}.}\label{fig:s2:capital-structure-arbitrage:curve}
\end{omfigure}

\section{The trade and its hedge ratio}

\begin{definition}[Cross-asset hedge ratio]\label{def:s2:capital-structure-arbitrage:hedge}
A \emph{cross-asset hedge ratio}\index{cross-asset hedge ratio} is the quantity of one asset that offsets the model sensitivity of a position in
another asset of the same issuer; in capital-structure arbitrage, the shares per unit of CDS notional that offset the CDS's value change when the
share price moves, minus the risky annuity times the implied spread's derivative in the share price.
\end{definition}

A trader who finds the market spread above the implied one sells protection: the CDS is paying too much for a risk the equity says is smaller. If the
share then falls, the implied spread rises and the protection seller loses. So the trader also shorts shares, as many as the hedge ratio says. At a
share price of 50 and debt of 40, that is 0.00067 shares per unit of notional: about 6,700 shares, or \$0.34 million of stock, against \$10 million of
protection. A buyer of protection does the reverse and buys the shares.

\texttt{firm.capstruct} runs the trade on sixty firms for ten years (\cref{lst:s2:capital-structure-arbitrage:trades}). Each firm's market spread is
the model's spread at the firm's true barrier and debt, times a mean-reverting mispricing (10\% in logs, half-life 34 days). The trader does not know
a firm's true barrier and assumes $L = 0.5$ everywhere, while the truth varies by firm with a dispersion of 25\%. The trader estimates the equity vol
from the trailing year. It enters when the market spread is more than 40\% (in logs) from the implied spread, and exits when the gap falls below 5\%
or after 180 days.

\section{Evidence and losses}

\begin{center}
\small
\begin{tabular}{@{}lccc@{}}
\toprule
sixty firms, ten years & leveraged buyouts 0.05 a firm-year & none & 0.2 a firm-year\\
\midrule
trades & 440 & 413 & 433\\
mean per trade (bp of notional) & 26.1 & 25.3 & 9.3\\
share losing & 48\% & 44\% & 52\%\\
worst trade (bp) & $-347$ & $-179$ & $-378$\\
share closed at 180 days & 56\% & 55\% & 57\%\\
Sharpe ratio of the book & 1.40 & 1.66 & 0.39\\
\bottomrule
\end{tabular}
\end{center}

The trades look like Yu's. Almost half lose, the worst lose several percent of notional, and the book of all of them earns a steady Sharpe ratio
(\cref{fig:s2:capital-structure-arbitrage:book}). Most of the return comes from the CDS leg, 20.4 of the 26.1 basis points, and the share hedge adds
5.6. More than half the trades never converge and are closed at 180 days. They were entered on a gap that was really the trader's wrong barrier, not a
mispricing, and those gaps do not close. Duarte, Longstaff and Yu found that the fixed-income arbitrage strategies needing more ``intellectual capital''
earned significant alphas after bond and equity factors and after fees, and that many had positively skewed returns, not the negative skew of
picking up nickels in front of a steamroller.

\begin{omfigure}
\begin{tikzpicture}
\begin{axis}[width=11cm,height=5.6cm,xlabel={\small year},ylabel={\small cumulative (\% of notional)},tick label style={font=\footnotesize},
  xmin=1,xmax=10,ymin=-1,ymax=5,grid=major,grid style={gray!25},table/col sep=comma,
  legend style={font=\scriptsize,at={(0.03,0.97)},anchor=north west}]
\addplot[omThm,thick] table[x=year,y=base]{figdata/strategies-2/09-capital-structure-arbitrage/book.csv};
\addplot[omDef,thick,dashed] table[x=year,y=shifts]{figdata/strategies-2/09-capital-structure-arbitrage/book.csv};
\legend{buyouts 0.05 a firm-year,buyouts 0.2 a firm-year}
\end{axis}
\end{tikzpicture}

\omcaption{The synthetic capital-structure book's cumulative P\&L per average open trade (about 27 at a time), as a share of one trade's notional,
with leveraged buyouts at 0.05 and 0.2 a firm-year. Data: \texttt{s2\_capstruct.cumulative}.}\label{fig:s2:capital-structure-arbitrage:book}
\end{omfigure}

\section{Model risk}

A leveraged buyout breaks the trade in both legs at once. The firm's debt doubles and its share jumps 20\% on the takeover premium, and the trader
learns the new debt from filings three months later. The market spread jumps at once. A trader who had sold protection loses on the CDS, and loses
again on the short shares, which rose. Of the trades open when their firm's leverage shifted, the five that had sold protection lost 98 basis points
on average (the worst 316). The seven that had bought protection earned 220. When buyouts come four times as often, 23 protection sellers were caught
and lost 174 basis points each on average.

The structural model is the trade's only measure of mispricing, and every input it takes is uncertain: the barrier, the equity vol, the debt a share.
A persistent gap is more often a wrong input than a market error. The trader needs rules for when to stop believing the model. Close gaps that do not
converge. Leave out firms in play for a takeover. Update debt from filings as they arrive.

\section{Strategy files}

\begin{strategyfile}{Equity-implied spread convergence}\label{strat:s2:capital-structure-arbitrage:convergence}
\sfield{Who pays you, and why} Investors in one market who price a company's default risk differently from investors in the other, until the two
converge.
\sfield{Instruments and venues} Single-name CDS; the company's shares for the hedge.
\sfield{Signal} Market spread against the \omterm{def:s2:capital-structure-arbitrage:implied}{equity-implied spread} from a structural model.
\sfield{Sizing and execution} Many small positions; hedge ratio from the model; exit on convergence or after a time limit.
\sfield{Costs} CDS bid-ask spreads; borrow on short shares.
\sfield{How it dies} Gaps that are model error; takeovers and recapitalisations that move equity and credit the same way.
\sfield{Horizon, capacity, infrastructure} Months; a structural model fed with filings, equity vol and CDS quotes.
\sfield{Backtest honestly} Debt as known at the time; include the trades closed at the time limit.
\sfield{Sources} Yu (2006); this chapter: a book Sharpe ratio of 1.40 and 48\% of trades losing.
\end{strategyfile}

\begin{strategyfile}{Long CDS, short equity on weak equity}\label{strat:s2:capital-structure-arbitrage:weak}
\sfield{Who pays you, and why} Credit investors slower than equity investors to price bad news.
\sfield{Instruments and venues} CDS protection bought; shares sold short.
\sfield{Signal} A falling share price with a CDS spread below its implied level.
\sfield{Sizing and execution} Sized to the hedge ratio; small in names close to distress.
\sfield{Costs} Protection premium; borrow on hard-to-borrow shares.
\sfield{How it dies} Recovery of the shares; borrow recall.
\sfield{Horizon, capacity, infrastructure} Weeks to months.
\sfield{Backtest honestly} Borrow costs on falling shares.
\sfield{Sources} No performance figure verified.
\end{strategyfile}

\begin{strategyfile}{Debt--equity dislocation after events}\label{strat:s2:capital-structure-arbitrage:events}
\sfield{Who pays you, and why} Holders forced to sell one side after an event (a downgrade, an index exclusion, a buyout announcement).
\sfield{Instruments and venues} CDS, bonds and shares of the same issuer.
\sfield{Signal} The event and the direction it moves each claim.
\sfield{Sizing and execution} Small; the event may change the capital structure itself.
\sfield{Costs} Wide spreads around events.
\sfield{How it dies} The event is a leverage shift: equity and credit are both repriced, correctly.
\sfield{Horizon, capacity, infrastructure} Days to weeks; event monitoring.
\sfield{Backtest honestly} Event dates and filings as known then.
\sfield{Sources} No performance figure verified; this chapter: protection sellers open through a buyout lost 98 basis points on average.
\end{strategyfile}

\section{Tutorial: same default, two prices}\label{tut:s2:capital-structure-arbitrage:tut}

\begin{tutorial}
\textbf{Goal.} Simulate firms with CDS and equity driven by a structural model with planted noise and leverage shifts, trade convergence, and measure
the losses when the model is wrong. \textbf{End state:} the table and the two figures.
\begin{tutsteps}
\item \textbf{The trades}.
\omcode{code/firm/capstruct/firm_capstruct.py}{118}{154}{Entry on the gap, CDS and share legs day by day, exit on convergence or time.}\label{lst:s2:capital-structure-arbitrage:trades}
\item \textbf{Results}.
\omcode{code/strategies-2/09-capital-structure-arbitrage/python/s2_capstruct.py}{33}{54}{Per-trade and book statistics, and trades caught by a leverage shift.}\label{lst:s2:capital-structure-arbitrage:stats}
\item \textbf{Run} \texttt{stats()}, \texttt{stats(0.2)}, \texttt{through\_shifts()}, \texttt{curve()} and \texttt{fig\_capstruct.py}.
\end{tutsteps}
\textbf{What to change next.} Give the trader the true barrier and measure how much of the loss was model error; skip firms whose share jumped more
than 15\% in a day; enter only after the gap has persisted for a month.
\end{tutorial}

\section{Build: capital-structure arbitrage}\label{bld:s2:capital-structure-arbitrage:capstruct}

\begin{build}
\bfield{Purpose} \omterm{def:s2:capital-structure-arbitrage:implied}{Equity-implied spreads} from a first-passage model, hedge ratios, convergence trades and a leverage-shift stress.
\bfield{Interface} \texttt{CapConfig(\dots)}, \texttt{implied\_spread(E, D, sigma\_E, cfg)}, \texttt{hedge\_ratio(E, D, sigma\_E, cfg)},
\texttt{simulate\_firms(cfg)}, \texttt{trades(sim, cfg)}.
\bfield{Rules} The trader sees debt with a lag and assumes one barrier for all firms; one unit of notional per trade; costs at entry and exit.
\bfield{Acceptance tests} \texttt{code/firm/capstruct/tests/}: spreads rise with leverage and fall with equity; the survival formula by hand; no
mispricing and no model error give no trades.
\bfield{Stretch} Barrier uncertainty as in CreditGrades; bond legs; recovery risk.
\end{build}

\begin{omsources}
\item F.~Yu, ``How profitable is capital structure arbitrage?'', \emph{Financial Analysts Journal} 62(5), 2006.
\item J.~Duarte, F.~A.~Longstaff and F.~Yu, ``Risk and return in fixed-income arbitrage: nickels in front of a steamroller?'', \emph{Review of
Financial Studies} 20(3), 2007.
\end{omsources}

\section{Exercises}

\begin{exercise}[$\star$]\label{exo:s2:capital-structure-arbitrage:1}
A five-year survival probability is 0.954 and recovery is 40\%. What spread does the chapter's formula give?
\end{exercise}

\begin{exercise}[$\star$]\label{exo:s2:capital-structure-arbitrage:2}
The hedge ratio is 0.00067 shares per unit of notional and the share price is 50. How much stock hedges \$10 million of protection?
\end{exercise}

\begin{exercise}[$\star$]\label{exo:s2:capital-structure-arbitrage:3}
A protection seller with a risky annuity of 4.5 sees the spread rise from 56 to 99 basis points. What does it lose per unit of notional?
\end{exercise}

\begin{exercise}[$\star\star$]\label{exo:s2:capital-structure-arbitrage:4}
Why do more than half the synthetic trades close at the time limit rather than on convergence?
\end{exercise}

\begin{exercise}[$\star\star$]\label{exo:s2:capital-structure-arbitrage:5}
Why does a leveraged buyout hurt a protection seller in both legs?
\end{exercise}

\begin{exercise}[$\star\star$]\label{exo:s2:capital-structure-arbitrage:6}
Yu found large losses on single trades and a steady portfolio. How can both be true?
\end{exercise}

\begin{exercise}[$\star\star\star$]\label{exo:s2:capital-structure-arbitrage:7}
\emph{Coding.} Run \texttt{stats(0.0)}. What does removing leverage shifts do to the worst trade and the Sharpe ratio, and what risk remains?
\end{exercise}

\begin{exercise}[$\star\star\star$]\label{exo:s2:capital-structure-arbitrage:8}
\emph{Find the flaw.} ``The CDS has been 50\% above its equity-implied level for six months, so the convergence trade is more certain than ever.''
\end{exercise}

\section{Problem: Same Default, Two Prices}

\begin{problem}[{Weekend problem --- capital-structure arbitrage}]\label{pb:s2:capital-structure-arbitrage:1}
The chapter's synthetic firms and the public record.

\medskip\noindent\textbf{Part I --- The model.}
\begin{enumerate}
\item Describe the first-passage model and its inputs.
\item Define the \omterm{def:s2:capital-structure-arbitrage:implied}{equity-implied spread} and give it at a share price of 50 with debt of 40 and 80.
\item Why do the model's spreads stay moderate as the share falls?
\item What did Yu use and find?
\end{enumerate}

\medskip\noindent\textbf{Part II --- The trade.}
\begin{enumerate}[resume]
\item Define the \omterm{def:s2:capital-structure-arbitrage:hedge}{cross-asset hedge ratio} and compute the hedge for \$10 million.
\item Describe the entry and exit rules.
\item What does the trader not know about each firm?
\item Where does most of the return come from?
\end{enumerate}

\medskip\noindent\textbf{Part III --- The results.}
\begin{enumerate}[resume]
\item Give the book's statistics with buyouts at 0.05 a firm-year.
\item Why do so many trades close at the time limit?
\item What did Duarte, Longstaff and Yu find?
\item How does the book compare with no buyouts?
\end{enumerate}

\medskip\noindent\textbf{Part IV --- The verdict.}
\begin{enumerate}[resume]
\item State the \emph{named result}: the convergence trade's Sharpe ratio and the losses when leverage shifts.
\item What happened to trades open through a buyout, by side?
\item Why is the barrier the model's weakest input?
\item What rules would you add?
\item How would you backtest the trade honestly?
\item Which strategy file is most exposed to buyouts?
\item How does this chapter relate to Book~6's structural model?
\item In one sentence: what does a capital-structure arbitrageur bet on?
\end{enumerate}
\end{problem}

\section{Interview questions}

\begin{interviewq}[$\star$ \iqroles{researcher}]\label{iq:s2:capital-structure-arbitrage:1}
How would you get a credit spread from a share price?
\end{interviewq}

\begin{interviewq}[$\star\star$ \iqroles{trader}]\label{iq:s2:capital-structure-arbitrage:2}
A company's CDS is far above its equity-implied level. Walk through the trade and what could go wrong.
\end{interviewq}

\begin{interviewq}[$\star\star$ \iqroles{risk}]\label{iq:s2:capital-structure-arbitrage:3}
How would you stress a capital-structure book?
\end{interviewq}

\begin{interviewq}[$\star\star$ \iqroles{researcher}]\label{iq:s2:capital-structure-arbitrage:4}
How would you tell a mispricing from a model error?
\end{interviewq}

\begin{interviewq}[$\star\star$ \iqroles{developer}]\label{iq:s2:capital-structure-arbitrage:5}
What data does the book need each day, and which of it arrives late?
\end{interviewq}

\begin{interviewq}[$\star\star\star$ \iqroles{researcher}]\label{iq:s2:capital-structure-arbitrage:6}
For assets following a driftless geometric Brownian motion (in log terms, drift $-\sigma^2/2$) and a constant barrier $B$, derive the probability of
not touching the barrier by time $t$ from the reflection principle.
\end{interviewq}
